\section{Passage from harmonic paths to primes}\label{sec:prime-paths}

We now transfer the path estimates to actual primes.  Throughout this
section, a \emph{standard prime path} starts at a lower state with cutoff
$B$ and gap $r_{\mathrm{st}}=s_{\mathrm{st}}B$, where
$s_{\mathrm{st}}\in[1.99,2.3]$.  Its primes decrease strictly after the
first step.  At the first step the upper endpoint is included; all
subsequent upper endpoints are excluded.  The next exponent is required
to exceed $\ell$, where $\ell\in\{1,2\}$.  A lower state includes every
available child, and an upper state includes precisely the children for
which $r-x\geq 2x$.  Thus these paths contain the paths admitted by
\eqref{eq:admission}, and agree with them above exponent $2$.

For a prefix $\mathfrak p=(p_1,\ldots,p_m)$ write
\[
 h(\mathfrak p)=\prod_{j=1}^m\frac1{p_j},\qquad
 r(\mathfrak p)=r_{\mathrm{st}}-\sum_{j=1}^m x(p_j),\qquad
 s(\mathfrak p)=\frac{r(\mathfrak p)}{x(p_m)}.
\]
The empty prefix has weight $1$.  Sums over prefixes below include all
allowed lengths, and therefore count a path once at each of its visited
states.  At a nonempty prefix, $x(p_m)>\ell$ is equivalent to
$s(\mathfrak p)<r(\mathfrak p)/\ell$.

We use the elementary consequences of \eqref{eq:parameters}
\begin{equation}\label{eq:prime-scales}
 \log w\sim\log B,\qquad
 S=(\log B)^2,\qquad N=O(S\log B)=O((\log B)^3).
\end{equation}
In particular, every fixed power of $S$ and $N$ is smaller than
$\exp(c\sqrt{\log w})$ for every fixed $c>0$, eventually.

\begin{lemma}[Harmonic prime measure]\label{lem:prime-harmonic}
There are absolute positive constants $C,c$ such that, on any exponent
interval with lower endpoint $u\geq1/2$,
\begin{equation}\label{eq:prime-harmonic}
 \sum_{x(p)\ \mathrm{in\ the\ interval}}\frac1p
   =\int_{\mathrm{interval}}\frac{dx}{x}+O(E(u)),
 \qquad E(u)=C\exp(-c\sqrt{u\log w}).
\end{equation}
The error is uniform in the upper endpoint and in either endpoint
convention.  In particular, the total weight of all decreasing prime
prefixes with exponents in $(\ell,B]$, with no admission restrictions,
is $O(B)$.
\end{lemma}

\begin{proof}
Apply partial summation to the prime number theorem in
Lemma~\ref{lem:prime-inputs}, between $w^u$ and the upper endpoint.
The main term is $\int dt/(t\log t)=\int dx/x$.  The integrated error,
including its lower-endpoint term, is bounded by a constant times
$\exp(-c\sqrt{\log(w^u)})$, after reducing $c$.  The bound remains valid
if the upper endpoint is sent to infinity in the error integral.  A
change of endpoint convention contributes at most one reciprocal prime,
which is absorbed by $E(u)$.  Finally,
\[
 \sum_{\mathfrak p}h(\mathfrak p)
   =\prod_{w^\ell<p\leq w^B}(1+1/p)
   \leq\exp\left(\sum_{w^\ell<p\leq w^B}\frac1p\right)\ll B.
\]
Deleting an endpoint prime or imposing admission restrictions can only
decrease this sum.
\end{proof}

\begin{lemma}[Removal of high states]\label{lem:high-states}
Fix $K>0$ and $A>0$.  The total original harmonic weight of prefixes
ending at gap at most $K$ and having some earlier or terminal ratio
greater than $S$ is $O_{K,A}(B^{-A})$.  More generally, for any fixed
$c'>0$, the sum of
$h(\mathfrak p)\exp(-c'r(\mathfrak p))$ over prefixes having such a
high state is $O_{A,c'}(B^{-A})$.  Both conclusions also hold in the
continuous harmonic model, uniformly in the starting ratios under
consideration.
\end{lemma}

\begin{proof}
Suppose a prefix reaches a state of gap $R$, cutoff $b$, and ratio
$R/b>S$.  Every subsequent factor has exponent at most $b$.  Reducing
the gap below $R/2$ therefore requires at least $S/2$ additional factors.
If $T=\log B+O(1)$ is an upper bound for the remaining harmonic
intensity, the sum of weights of these continuations is at most
\[
 \sum_{j\geq S/2}\frac{T^j}{j!}.
\]
Since $S=(\log B)^2$, this tail is smaller than every fixed negative
power of $B$, even after multiplication by $B$.  To sum over all
possible high ancestors, use Lemma~\ref{lem:prime-harmonic}: their
total prefix weight is $O(B)$.  This argument is an upper bound even
when it counts a terminal prefix more than once.

At a high ancestor $R>Sb>S\ell$.  Thus a prefix ending at fixed gap
$K$ must reduce this gap by more than half for all sufficiently large
$B$.  For the weighted assertion, prefixes which do not do so have
terminal gap greater than $S\ell/2$, and their reward is
$O(\exp(-c'S\ell/2))$.  Their total unweighted prefix mass is $O(B)$,
so they too contribute less than every $B^{-A}$.  In the continuous
model, the ordered integrals of length $j$ are bounded by $T^j/j!$,
and their total is $e^T\ll B$; the proof is unchanged.
\end{proof}

For the rest of the discrete approximation, retain only states of ratio
at most $S$.  There are at most $N=C S\log B$ arrived-at states, with
a fixed sufficiently large $C$.  Indeed a transition to ratio $t\leq S$
decreases $\log r$ by $g(t)=\log(1+1/t)\gg1/S$.  Moreover the gap at a
retained nonempty state is bounded below by $0.99\ell$, and the starting
gap is $O(B)$.  This also proves the same deterministic horizon for a
continuous path up to its first exit from these conditions.

\begin{proposition}[Coupling the tilted paths]\label{prop:prime-coupling}
Choose sufficiently small absolute positive constants in the following
order.  For a permitted prime child of a state of type $i$, ratio $s$,
and gap $r$, put $t=r/x(p)-1$ and give that child mass
\begin{equation}\label{eq:prime-tilt}
 q(p\mid r,s,i)=\frac1p\left(\frac{t+1}{t}\right)^2
                    \frac{\phi_{i'}(t)}{\phi_i(s)}.
\end{equation}
Only children with $x(p)>\ell$ and $t\leq S$ are retained.  For
$\epsilon_w=\exp(-c_2\sqrt{\log w})$ their total mass is at most
$1+\epsilon_w$.  Dividing by $1+\epsilon_w$ and assigning the remaining
mass to a cemetery state defines a killed Markov chain.

This chain can be coupled with the continuous chain
\eqref{eq:tilted-kernel}, killed at the first draw with $x\leq\ell$ or
$t>S$, so that the probability of failure is at most
$\exp(-c_3\sqrt{\log w})$.  On success the chains have the same number
of arrived-at states and matching exits.  At corresponding states,
\begin{equation}\label{eq:prime-coupling-errors}
 |s-\widetilde s|\leq h,\qquad
 \left|\log\frac r{\widetilde r}\right|\leq C N h,\qquad
 h=\frac1{\lceil\exp(c_4\sqrt{\log w})\rceil}.
\end{equation}
All statements are uniform in $s_{\mathrm{st}}\in[1.99,2.3]$ and in the
endpoint conventions stated above.  For a length-$m$ retained prefix,
if $\mathbb P_w(\mathfrak p)$ denotes its probability in the killed
chain, then the exact weight identity is
\begin{equation}\label{eq:prime-weight}
 h(\mathfrak p)=\mathbb P_w(\mathfrak p)(1+\epsilon_w)^m
 \left(\frac{r(\mathfrak p)}{r_{\mathrm{st}}}\right)^2
 \frac{\phi_\ee(s_{\mathrm{st}})}{\phi_i(s(\mathfrak p))},
 \qquad (1+\epsilon_w)^m=1+o(1).
\end{equation}
\end{proposition}

\begin{proof}
Under $t=r/x-1$ we have $|dx/x|=dt/(t+1)$.  Multiplying by the
non-prime part of \eqref{eq:prime-tilt} gives
$W_0(t)\phi_{i'}(t)/\phi_i(s)$, exactly the continuous transition
density.  Also $r_{\mathrm{child}}/r=t/(t+1)$, so multiplication over a
prefix telescopes and proves \eqref{eq:prime-weight}.  The last
assertion there follows from $m\leq N$ and \eqref{eq:prime-scales}.

Partition the ratio axis into bins $[jh,(j+1)h)$.  For a bin meeting
the allowed next range, temporarily extend it to the full bin when
seeking an upper bound.  Its corresponding $x$-interval has logarithmic
width comparable to $h/(t+1)$ and lower endpoint at least $\ell/2$,
for sufficiently large $w$.  Lemma~\ref{lem:derivative-weights}
bounds logarithmic derivatives of the weights and backward opposite
weight ratios by fixed polynomials in $S$.  Therefore a full interior
bin has discrete mass equal to its continuous mass with relative error
\begin{equation}\label{eq:prime-bin-error}
 O\left(\operatorname{poly}(S)
             \left(h+\frac{E(1/2)}h\right)\right).
\end{equation}
Here the first term controls variation across the bin, and the second
is \eqref{eq:prime-harmonic} divided by its harmonic main mass.
The density itself is bounded by $\operatorname{poly}(S)$ on any
bin meeting the support.  There are only a bounded number of partial
support bins, each with mass at most
$\operatorname{poly}(S)h$.  The fixed even-child boundary $2$ is a
mesh endpoint.  Summing the interior relative errors and the partial
bin upper bounds proves that the total mass is at most
$1+\epsilon_w$, after choosing $c_4$ and then $c_2$ sufficiently small.

Suppose the chains have been matched through the previous draw.  Their
lower support endpoints are $s-1$ or $\max(2,s-1)$, and consequently
differ by at most $h$.  The upper endpoint imposed by $x>\ell$ is
$r/\ell-1$.  It needs comparison only when it is at most $S$, up to a
mesh error; in that case $r=O(S+1)$, and its discrepancy is at most
$O((S+1)Nh)$.  Thus a large starting gap does not enter this support
error.  Ratios of the two density denominators are
$1+O(\operatorname{poly}(S)h)$.  The total mass in mismatched or partial
bins, together with the total variation distance between full bin
labels and the discrepancy in exit probabilities after normalization,
is at most $\exp(-c'\sqrt{\log w})$.

Maximally couple these bin labels and the cemetery label at each step.
Inside a matched nonexit bin the two next ratios differ by at most $h$.
Since $g'(t)=-1/(t(t+1))$ is bounded on the used ranges, each matched
step changes the logarithmic-gap discrepancy by $O(h)$.  Iteration up
to $N$ steps gives \eqref{eq:prime-coupling-errors}; a union bound and
\eqref{eq:prime-scales} give the stated failure probability.  Single
endpoint primes have already been included in the uniform estimate
\eqref{eq:prime-harmonic}.  The initial lower ratios slightly below $2$
are covered by the extended derivative-weight identities in
Lemma~\ref{lem:derivative-weights}.
\end{proof}

The total variation coupling by itself is insufficient for rewards
containing $1/\phi_i(s)$ far out in the ratio tail.  The following
pointwise majorants supply the required additional control.

\begin{lemma}[Discrete marginal majorants]\label{lem:discrete-marginals}
For the killed chain of Proposition~\ref{prop:prime-coupling}, the
probability of a nonexit bin with left endpoint $t$ at step $2\leq n\leq N$
is bounded, with an absolute constant, by
\begin{equation}\label{eq:prime-marginals}
 \begin{cases}
  C h\phi_\ee(t),&n\ \hbox{even},\\[2mm]
  C h\phi_\oo(t)\bigl(1-t^{-2}+3hW_0(t)\bigr),&n\ \hbox{odd}.
 \end{cases}
\end{equation}
At step $1$ the bound is $C hW_0(t)\phi_\oo(t)$, on the slightly
extended range $t\geq0.98$.  The constants are uniform in the starting
ratios and in $w$ sufficiently large.
\end{lemma}

\begin{proof}
The bin argument in the preceding proof gives the conditional upper
bound
\begin{equation}\label{eq:prime-conditional-bin}
 (1+\delta_w)hW_0(t)\frac{\phi_{i'}(t)}{\phi_i(s)},
 \qquad \delta_w=\exp(-c'\sqrt{\log w}),
\end{equation}
for every feasible child bin.  Division by $1+\epsilon_w$ cannot
increase it.  The first-step assertion follows because the starting
denominator is bounded away from zero.  At the second step, canceling
the intermediate derivative weight leaves an incoming sum bounded by
\[
 C\sum_{s_{\mathrm{st}}-1\leq u\leq t+1+h}hW_0(u)
       \ll \log(t+2).
\]
Since $t\geq2$ and $W_0(t)\log(t+2)$ is bounded, this proves the even
majorant at step $2$.

For the induction, replacing a weight at an actual state by its value
at the bin's left endpoint costs a relative
$1+O(\operatorname{poly}(S)h)$.  Given the even majorant, the incoming
bins for an odd child of left endpoint $t$ start at $2$ and end at most
at $t+1+h$.  Their total mesh length is at most $t-1+3h$.  Multiplication
by $W_0(t)$ gives
\[
 W_0(t)(t-1+3h)=1-t^{-2}+3hW_0(t),
\]
as required.  In the reverse direction, the incoming odd bins start at
$1$; their bracket sums satisfy
\begin{align*}
 \sum_{1\leq u\leq t+1+h}h
       \bigl(1-u^{-2}+3hW_0(u)\bigr)
 &\leq\int_1^{t+1}(1-u^{-2})\,du+O(h\log(t+2))\\
 &=\frac{t^2}{t+1}+O(h\log(t+2)).
\end{align*}
The new factor $W_0(t)$ cancels the leading expression, and the error
is $O(h)$ for $t\geq2$.  The points $1$ and $2$ are mesh aligned, and
the $3hW_0$ term covers the endpoint mesh interval where the stationary
odd factor vanishes.  Each iteration enlarges the constant by at most
$1+O(\delta_w+\operatorname{poly}(S)h)$.  Its product over $N$ steps is
$1+o(1)$ by \eqref{eq:prime-scales}, completing the proof.
\end{proof}

\begin{proposition}[Prime occupation estimates]\label{prop:prime-occupation}
The following assertions hold uniformly for
$s_{\mathrm{st}}\in[1.99,2.3]$ and $\ell\in\{1,2\}$.
\begin{enumerate}[label=\textup{(\roman*)}]
\item For fixed $K$, the original harmonic mass of prefixes with
terminal gap at most $K$ is $O_K(B^{-2})$.
\item Consider any selection of such compact-ending prefixes whose
low-state histories imply an event $\mathcal A$ in the killed prime
chain.  Their total weight is at most
\begin{equation}\label{eq:prime-compact-event}
 C_K B^{-2}\mathbb P_w(\mathcal A)+O_{K,A}(B^{-A})
\end{equation}
for every fixed $A>0$.  One may take $\mathcal A$ to be the event that
at least one selected compact prefix occurs.
\item For every fixed $c'>0$,
\begin{equation}\label{eq:prime-tail}
 \lim_{K\to\infty}\limsup_{w\to\infty}
 B^2\sum_{\substack{\mathfrak p\ \mathrm{standard}\\r(\mathfrak p)>K}}
       h(\mathfrak p)e^{-c'r(\mathfrak p)}=0.
\end{equation}
The same compact and tail assertions hold for continuous original
harmonic weights.
\item Let $H_i(r,s)$ be bounded continuous functions supported in a
fixed compact positive gap interval and a fixed compact ratio interval.
Fixed interval boundaries and the indicator $s<r/\ell$ are also
permitted, with one-sided continuous extension before imposing that
indicator, as in Proposition~\ref{prop:occupation}.  Then
\begin{align}\label{eq:prime-occupation}
 &\frac{r_{\mathrm{st}}^2}{\phi_\ee(s_{\mathrm{st}})}
   \sum_{\mathfrak p}h(\mathfrak p)
        H_{i(\mathfrak p)}(r(\mathfrak p),s(\mathfrak p))\notag\\
 &\hspace{8mm}\longrightarrow
 \frac1{M_g}\sum_{i\in\{\ee,\oo\}}
   \int_0^\infty\int
    \frac{r^2\pi_i(s)}{\phi_i(s)}H_i(r,s)
       \1_{s<r/\ell}\,ds\,\frac{dr}{r},
\end{align}
where the inner domains are $s\geq2$ for $i=\ee$ and $s\geq1$ for
$i=\oo$.  The convergence is uniform in the starting ratios.
\end{enumerate}
\end{proposition}

\begin{proof}
Remove high-state histories by Lemma~\ref{lem:high-states}.  On a
remaining compact-ending prefix, $s<r/\ell$ bounds the terminal ratio
in a fixed compact set.  The positive weights $\phi_i$ are bounded
away from zero there, so \eqref{eq:prime-weight} is at most
$C_K B^{-2}\mathbb P_w(\mathfrak p)$.  A path has at most $C_K$ visits
with gap at most $K$, since each exponent exceeds $\ell$ and thus each
step decreases the gap by more than $\ell$.  Summing proves (i) and
\eqref{eq:prime-compact-event}.  This also explains why a history
exception costs its tilted probability times $B^{-2}$, rather than
its probability times the unrestricted original mass $O(B)$.

For (iii), the scaled reward after the change of measure is at most
\begin{equation}\label{eq:prime-scaled-reward}
 C\frac{r^2e^{-c'r}}{\phi_i(s)}.
\end{equation}
Put $K_B=C'\log\log B$.  At a step covered by
Lemma~\ref{lem:discrete-marginals}, sum first over ratio bins.  Since
$r>\ell s$ and $\ell\geq1$, cancellation of $\phi_i$ bounds the
expected reward from $r>K_B$ by a constant times
\begin{equation}\label{eq:prime-far-tail}
 \sum_{t=jh\geq0.97}h
       \sup_{u\geq\max(t,K_B)}u^2e^{-c'u}
 \ll (1+K_B)^3e^{-c'K_B}.
\end{equation}
Increasing $C'$ makes this $o(N^{-1})$.  The first-step majorant gives
the same bound, because $W_0$ is bounded on its extended range.  The
empty prefix contributes $O(B^2e^{-c''B})$.

On $K<r\leq K_B$, the Dickman lower bounds in
Lemma~\ref{lem:derivative-weights} bound
\eqref{eq:prime-scaled-reward} by
$\exp(O(K_B\log(K_B+2)))$.  Even after multiplication by $N$, the
coupling failure contribution is negligible, since
\[
 \log N+O(K_B\log(K_B+2))=o(\sqrt{\log w}).
\]
On a successful coupling, the reward is bounded by
\[
 C\frac{\widetilde r^2e^{-c'\widetilde r/2}}
          {\phi_i(\widetilde s)},\qquad
 \widetilde r>K/2,\quad \widetilde s<\widetilde r/\ell.
\]
Here the indicator follows from matching the nonexit arrivals, and
the comparison of weights follows from their logarithmic derivative
bounds and \eqref{eq:prime-coupling-errors}.  On a band
$R\leq\widetilde r\leq eR$, Lemma~\ref{lem:band-bounds} bounds the
expected sum after cancellation of $\phi_i$ by $O(R+1)$.  Thus this
band contributes at most $C(R+1)R^2e^{-c'R/2}$.  Summing the geometric
sequence of bands above $K/2$ proves \eqref{eq:prime-tail}.  High-state
histories were already negligible by Lemma~\ref{lem:high-states}.
The continuous proof uses the same band bound directly, along with
the continuous version of that lemma.

Finally, for (iv) use \eqref{eq:prime-weight} with the compact test
$r^2H_i(r,s)/\phi_i(s)$.  The correcting factor is uniformly $1+o(1)$,
and the number of compact-gap visits is bounded.  Coupling therefore
transfers this test by uniform continuity.  Boundary indicators cause
no difficulty: the limiting continuous occupation measure gives zero
mass to their boundaries, and continuous upper and lower
approximations suffice.  Continuous removal of high states is
harmless as well.  For example, the drift estimate in
Lemma~\ref{lem:regeneration} gives a uniformly bounded exponential
moment of the ratio at each step, so a union bound up to the low-state
horizon gives $O(Ne^{-S})=o(1)$ for reaching a high state before exit.
Alternatively one can use the continuous original-weight conclusion
of Lemma~\ref{lem:high-states} on compacts.  Apply
Proposition~\ref{prop:occupation} to the transferred test and use its
uniformity in the starting state.  This gives
\eqref{eq:prime-occupation} and completes the proof.
\end{proof}
